\documentclass[10pt]{article}
\usepackage[margin=0.72in]{geometry}
\usepackage{amsmath,amssymb,amsthm,mathtools}
\usepackage[T1]{fontenc}
\usepackage{lmodern}
\usepackage{hyperref}
\usepackage{microtype}
\hypersetup{colorlinks=true,urlcolor=blue,linkcolor=blue}
\setlength{\parindent}{0pt}
\setlength{\parskip}{0.45em}
\newtheorem{theorem}{Theorem}
\newcommand{\Q}{\mathbb Q}
\newcommand{\Ext}{\operatorname{Ext}}
\title{Problem 00000000561: a square-zero Koszul algebra\protect\\ excludes every finite cyclotomic denominator}
\author{lizaixi01}
\date{24 October 2026}
\begin{document}
\maketitle
\vspace{-1.8em}

\textbf{Precise scope.}
The printed conjecture asks that the Hilbert series of the Yoneda algebra of an
$N$-Koszul algebra be rational with denominator $\prod_i(1-t^{d_i})$.
The product has no printed index set. We use the ordinary rational-function
meaning of a denominator: a \emph{finite polynomial} product, with positive
integer $d_i$, allowing repetitions. Under this reading the denominator clause
is false already for $N=2$. The example's series is itself rational.
We do not rule out an infinite formal Euler product, or dispute rationality for
this example. The full original English and Chinese text is in \texttt{SOURCE.md}.

\begin{theorem}
There is a connected quadratic $2$-Koszul $\Q$-algebra $A$, with augmentation
$k=\Q$, such that its actual graded Yoneda spaces satisfy
\[
 \dim_\Q\Ext_A^n(k,k)=2^n\quad(n\geq0),\qquad
 H(t)=\sum_{n\geq0}\dim_\Q\Ext_A^n(k,k)t^n=\frac1{1-2t}.
\]
For every finite family of positive integers $(d_i)$ and every $P\in\Q[t]$,
\[
 \left(\prod_i(1-t^{d_i})\right)H(t)\ne P(t).
\]
The empty family is allowed as well.
\end{theorem}

\section*{1. The algebra and its full linear minimal resolution}
Use the native square-zero extension
\[
 A=\Q\oplus V,\quad V=\Q^{\{0,1\}},\qquad
 (r,v)(s,w)=(rs,rw+sv).
\]
Let $\epsilon_a=(0,e_a)$ for the two genuine coordinate basis vectors $e_a$.
All four products $\epsilon_a\epsilon_b$ vanish. The grading is
$A_0=\Q\oplus0$, $A_1=0\oplus V$, and $A_j=0$ for $j\geq2$.
The formalization constructs its actual submodules, native internal direct-sum
decomposition and native \texttt{GradedAlgebra} structure. The degree-zero
component is linearly isomorphic to $\Q$ and the two $\epsilon_a$ span degree one.

This is a genuine quadratic presentation, rather than an assigned label:
\[
 A\simeq \Q\langle x_0,x_1\rangle/(x_ax_b\mid a,b\in\{0,1\}).
\]
Lean constructs the free algebra, the relation identifying each displayed
length-two monomial with zero, and its native \texttt{RingQuot}. The free
evaluation $x_a\mapsto\epsilon_a$ descends. In the reverse direction, a vector
$v\in V$ maps to $\sum_a v_a\bar x_a$; products of any two such images vanish.
The native square-zero-extension lift supplies the inverse homomorphism.
The two composites are proved to be identities on all elements.

Let $B_0$ be the one-element set and $B_{n+1}=\{0,1\}\times B_n$.
These are actual finite binary words, with cardinality proved inductively to
be $2^n$. Put $P_n=A^{B_n}$ and define actual $A$-linear maps
\[
 d_n:P_{n+1}\longrightarrow P_n,\qquad
 (d_nf)(w)=\sum_{a=0}^1\epsilon_a f(a,w),
 \qquad \eta:P_0\longrightarrow k,\quad\eta(f)=f(*)_0.
\]
Here $(r,v)_0=r$. Coordinate calculations give, for every $n,f,w,a$,
\[
 (d_nf)(w)_0=0,\qquad (d_nf)(w)_1(a)=f(a,w)_0.
\]
Consequently $d_nd_{n+1}=0$,
$\ker d_n$ consists exactly of functions with zero scalar coordinates,
and $\operatorname{im}d_{n+1}$ is exactly that same submodule.
For surjectivity onto this submodule, lift $g$ by
$f(a,w)=(g(w)_1(a),0)$. The same calculation proves
$\ker\eta=\operatorname{im}d_0$ and $\eta$ is surjective.
Thus the entire infinite complex is exact in every required degree.

Each $P_n$ has the native finite free coordinate basis and is projective.
Lean packages this exact augmented complex as an actual native
\texttt{ProjectiveResolution}, including its quasi-isomorphism to $k[0]$.
Give its basis vectors degree $n$: explicitly
$(P_n)_j=A_{j-n}^{B_n}$ for $j\geq n$, and zero otherwise.
The native component decompositions and graded scalar actions are constructed;
every $d_n$ preserves internal degree, and the augmentation preserves the
grading of $k$ concentrated in degree zero. On basis vectors,
\[
 d_n(e_{(a,w)})=\epsilon_a e_w.
\]
All matrix entries lie in degree one and in the actual augmentation ideal.
The image is proved to equal the $A$-span of these positive-degree basis
multiples. This is a minimal linear graded free resolution. It proves the
ordinary Koszul defining criterion \cite[Definition 8.1]{MP}; in the $N$-Koszul
notation $\nu_N(i)=N\lfloor i/2\rfloor+(i\bmod2)$, Lean proves
$\nu_2(i)=i$ for every $i$. This is the $N=2$ case of the standard convention
in \cite[p.~4, Definition 2.2]{Berger}. No theorem about Hilbert series is assumed.

\section*{2. Native Ext and the denominator obstruction}
For every actual $A$-linear map $P_n\to k$, precomposition with $d_n$ is zero,
because every $\epsilon_a$ acts as zero on $k$. Evaluation on the genuine free
basis gives a proved $\Q$-linear equivalence
\[
 \operatorname{Hom}_A(P_n,k)\simeq_\Q\Q^{B_n}.
\]
The inverse sends a coordinate family $c$ to
$f\mapsto\sum_w f(w)_0c_w$. Reconstruction is proved on every $f$.
The native projective-resolution comparison constructs all elements of the
native $\Ext_A^n(k,k)$ from such cocycles. Its actual surjectivity theorem and
coboundary criterion show that the map is bijective for $n>0$, since every
possible boundary map is zero. Compatibility with the native $\Q$-scalar
action is proved using Yoneda postcomposition. For $n=0$, native
$\Ext^0\simeq\operatorname{Hom}_A(k,k)$ followed by evaluation at $1$ gives
the actual equivalence with $\Q$.

All native Ext spaces are therefore proved finite dimensional, and their
finranks are derived from $|B_n|=2^n$. The Hilbert series is the actual formal
power series whose coefficient is this native finrank. In particular, no
infinite-dimensional finrank default is used. The proof uses the graded
native Ext spaces; it includes native bilinear Yoneda composition but does
not construct a separate bundled total Yoneda ring. Multiplication is not
needed for this Hilbert-series assertion.

The coefficient at zero is $1$ and successive coefficients double. Lean
therefore proves the formal-series identity
\[
 (1-2t)H(t)=1.
\]
Suppose a finite positive family gave $D(t)H(t)=P(t)$, where
$D(t)=\prod_i(1-t^{d_i})$. Multiplying by $1-2t$, commuting factors and using
this identity gives $(1-2t)P(t)=D(t)$ as an identity of actual formal series.
Injectivity of the native polynomial-to-power-series map makes this an
identity of polynomials. Evaluating at $t=1/2$ gives
\[
 0=D(1/2)=\prod_i(1-2^{-d_i})\ne0,
\]
because each $d_i>0$. Lean proves this nonvanishing for every finite index
type, including repeated exponents and the empty product. This contradicts
the finite-denominator clause independently of the source's particular
degree-sequence choice.

\newpage
\section*{3. Formalization boundary and reproduction}
All algebra, grading, full-degree exactness, projectivity, actual native Ext
comparisons, dimensions, formal-series operations and the finite-denominator
contradiction above are proved in \texttt{lean/Main.lean}. The conventional
name ``$2$-Koszul'' is identified by its defining minimal-linear-resolution
criterion; the proof supplies that criterion's full mathematical data.
The general terminology and free quadratic presentation convention come from
the cited definitions, not additional Lean axioms. In particular,
\cite[Theorem 8.5]{MP} is background only; its Poincare-series formula is not
used as an assumed computation. Selected complete pages 24--27 of that survey
and page 4 of Berger were inspected, not the complete papers.

Key declarations in namespace \texttt{TLMC561} are:
\begin{center}
\begin{tabular}{@{}l@{}}
\path{actualQuadraticPresentation}\\
\path{resolution}\\
\path{full_linear_resolution_certificate}\\
\path{graded_minimal_augmentation_certificate}\\
\path{ext_finrank}\\
\path{hilbert_identity}\\
\path{no_finite_cyclotomic_denominator}
\end{tabular}
\end{center}
The exhaustive check enumerates the actual
namespace, including generated declarations, and issues ordinary
\texttt{\#print axioms} for every name. Only \texttt{propext},
\texttt{Classical.choice} and \texttt{Quot.sound} are allowed; there is no
\texttt{sorry}, \texttt{native\_decide} or additional axiom.

The project pins Lean 4.33.0 and Mathlib commit
\path{db584cd6d46c92f209a44c0f1c829460d327499d}.
From \texttt{lean/}, run \texttt{lean --version}, then \texttt{./verify.sh}
for the build, warning-as-error axioms audit and official fresh-kernel replay.
Equivalently run the latter three commands separately:
\begin{verbatim}
lake build +Main
lake env lean -DwarningAsError=true Check.lean
lake env leanchecker --fresh Main
\end{verbatim}
The evidence directory records source-bound commands, logs and observed
thread affinity; production timings are diagnostic, not a controlled
research-efficiency result. Organizer acceptance is pending.

\begin{thebibliography}{9}
\bibitem{MP} J. McCullough and I. Peeva,
\emph{Infinite Graded Free Resolutions}, August 2018 survey,
pp.~24--27, especially Definition 8.1 and the discussion of Yoneda algebras.
\url{https://pi.math.cornell.edu/~irena/papers/infiniteRes.pdf}.
\bibitem{Berger} R. Berger,
\emph{Gerasimov's theorem and $N$-Koszul algebras},
arXiv:0801.3383v3, 19 December 2008, p.~4, Definition 2.2.
\url{https://arxiv.org/abs/0801.3383}.
\end{thebibliography}
\end{document}
